\section{PDDL：Neuro-Symbolic Planning}

\subsection{语言模型直接规划的困境}

在 Blocked World 等经典 Planning 问题中（如积木世界），语言模型直接产生 Plan 时很容易出错。例如 GPT-4 在第四步就产生了错误——因为涉及状态约束（机械臂一次只能拿一个 block、block 必须在桌面或某一个位置上），语言模型很容易在内部产生明显的逻辑错误。

\subsection{LM+P：扬长避短的方案}

\begin{importantbox}{LM+P 的核心思想}
\begin{itemize}
    \item \textbf{语言模型不擅长的}：执行精确的 Reasoning 或解决 Planning 问题
    \item \textbf{语言模型擅长的}：翻译——将自然语言转换为符号语言
    \item \textbf{PDDL Solver 擅长的}：基于完整的 Domain 定义和 Goal State，产生可靠的执行序列
\end{itemize}
因此，让语言模型做翻译，让 PDDL Solver 做规划求解。
\end{importantbox}

具体流程：
\begin{enumerate}
    \item 有一个 Domain 的 PDDL 定义（描述环境中有哪些 action、precondition、effect）
    \item 用户给出自然语言指令（如"把苹果放在冰箱里"）
    \item 语言模型\textbf{翻译}为 Goal State（如 apple\_1 在 fridge\_1 中）
    \item PDDL Solver 基于 Domain 定义和 Goal State，\textbf{求解}出具体的 Plan
    \item 语言模型再把每个步骤\textbf{翻译}回自然语言
\end{enumerate}

\begin{warningbox}{语言模型不应被寄予过高期望}
即使是强如 Gemini 2.5 Pro，在解方程 $5.9 = x + 5.11$ 这样简单的问题时也会犯错。我们不应该预期语言模型可以 solve 所有问题——这时一个简单的计算器 tool 就够了。要做系统级的设计，在合适的环节使用合适的工具。
\end{warningbox}

\subsection{本章小结}

PDDL 代表了 Neuro-Symbolic 方案的精髓：让语言模型做它擅长的翻译，让专业的符号求解器做精确的规划。这种扬长避短的设计思路是构建可靠 Agent 的关键。


